%Do not modify this content without consulting with its authors!!!
%Jakub & Olek

\subsection{Future Research Directions}
Although GenAI solutions have seen impressive development in recent years, progress in this field continues to raise new questions. This is especially true when considering a wider context beyond a technical focus. Successful application of GenAI systems and related digital transformations requires addressing many challenges and providing solutions that are much more mature than those currently available. This creates the need for future research in many areas and directions. 

\subsubsection{User Perspective [F1]}

\paragraph{User-GenAI Interaction Design}
Future research on user-GenAI interaction design in IS should prioritize the development of trustworthy, transparent, and explainable AI systems, as these themes recur across multiple sources \parencite{R02, R08, S18}. In order to foster user's trust it is necessary to address issues such as hallucinations, inherited biases, and the interpretability of LLM outputs. Another identified direction is the need for personalization and adaptive design, which includes tailoring GenAI systems to individual cognitive and emotional needs, supporting diverse user routines (both professional and private), and enabling proactive emotional management \parencite{R05, R08, S11}. Cultural and linguistic sensitivity is another prominent direction, with calls to localize GenAI systems and mitigate stereotypes in interactions \parencite{R01, S11}. The integration of multimodal capabilities and sensory engagement further expands the design space for user experience \parencite{S01, R08}. Finally, IS researchers are encouraged to explore the broader systemic implications of GenAI, including its impact on digital work, human-machine symbiosis, and the evolving boundaries of end-user computing \parencite{R04, R07, R10, S08}.

\paragraph{The Required Skills and Education}
The rise of GenAI is prompting a fundamental rethinking of the skills and educational frameworks required across industries. A recurring theme in the literature is the need to define and cultivate new competencies, including AI literacy and ethical usage, especially among non-technical users \parencite{R07, S06, S12}. Researchers are encouraged to explore how training programs and professional development initiatives can equip both workers and educators to effectively integrate GenAI into their workflows and pedagogical practices \parencite{R08, S12}. Closely linked is the transformation of job roles and labor markets, with GenAI expected to automate routine tasks, reshape existing roles, and create entirely new categories of work \parencite{R04, R05, R07, R08}. Understanding which tasks are most affected and how workers can adapt is a critical area for future inquiry. Additionally, the integration of GenAI into organizational contexts—such as IT outsourcing, service marketing, and mission-critical domains like healthcare and finance—will require a reassessment of workforce capabilities and strategic positioning \parencite{R03, R10}. Finally, the emergence of AI agents as team members introduces new questions about the competencies needed to work effectively alongside intelligent systems \parencite{S06}.

\paragraph{Personalization Tailored to the Individual User}
Personalization in GenAI systems is an emerging topic in IS research, with multiple sources emphasizing the need to tailor interactions to individual users’ preferences, values, and backgrounds \parencite{R05, R08, S18}. Future studies should explore how AI-augmented services can proactively support users through personalized prompts, assistance, and emotional alignment, while also adapting to diverse cognitive and behavioral patterns. Closely related is the challenge of balancing hyper-personalization with ethical accountability, particularly in marketing and customer engagement contexts where biases may be amplified \parencite{R01, R08}. Another important direction involves understanding how GenAI affects individual users—both workers and consumers—and how these technologies can enhance satisfaction, engagement, and perceived significance in digital interactions \parencite{R07, R10}. Finally, personalization of explanations and communication, especially in domains like healthcare and education, is seen as a key factor for improving user comprehension and decision-making \parencite{S11, S18}.

\paragraph{Psychological Effects of GenAI Use on Individuals}
The psychological effects of GenAI on individuals represent a multifaceted research area within IS domain. A dominant theme across sources is the need to understand how continuous interaction with GenAI influences users’ cognitive processes, emotional states, and work habits over time \parencite{R05, R07, R08, R09, R10}. Researchers are encouraged to investigate both the risks of over-reliance and the potential for loss of control due to over-automation, as well as strategies to foster healthy human-AI relationships \parencite{R07, R08}. Closely related is the impact of GenAI on decision-making and information-seeking behavior, which may alter users’ autonomy \parencite{R07}. The evolving dynamics of social and personal interactions mediated by GenAI also call for interdisciplinary, user-centered approaches that account for ethical and emotional dimensions \parencite{R07, R09, S08}. Moreover, future studies should explore how GenAI systems can be designed to proactively manage emotional experiences, especially during periods of uncertainty and change \parencite{R08}. Overall, this research agenda highlights the importance of capturing the nuanced and long-term psychological implications of GenAI use across diverse user populations.

\paragraph{Risks Affecting Individual Users}
The increasing integration of GenAI into everyday digital experiences raises significant risks for individual users, prompting a need for focused IS research. A recurring concern is user's privacy and data security, especially in real-world and real-time chatbot interactions where user data may be exposed without sufficient safeguards \parencite{S08, S11}. Researchers are urged to develop ethical and regulatory frameworks to address these vulnerabilities. Additionally, the amplification of cyber threats through GenAI and the spread of AI-generated disinformation call for advanced countermeasures and real-time detection systems \parencite{R03, R04}. Another identified risk concerns over-reliance on GenAI and the psychological consequences of excessive automation, including loss of control and technological resistance \parencite{R07, S12}. Understanding these risks is essential to ensure safe and responsible use of GenAI technologies.

\subsubsection{Organizational Perspective [F2]}

\paragraph{GenAI Adoption and Acceptance Factors}
A prominent direction for future research on GenAI adoption and acceptance in IS involves understanding the trust-related factors that influence user interaction with LLMs. Multiple sources emphasize the importance of explainability, interpretability, and transparency as key enablers of trust and acceptance, particularly in  domains such as healthcare and education \parencite{R02, R09, S01, S07}. These factors are closely tied to concerns about bias, hallucinations, and data quality, which design science researchers are encouraged to address through novel system architectures and validation methods \parencite{R02, S07}. Another widely discussed area is the organizational and cultural context of GenAI adoption, including how organizational norms, structures, and processes shape individual and collective attitudes toward automation and AI integration \parencite{R05, R07, R10}. Additionally, scholars call for investigations into sector-specific adoption dynamics, such as in education \parencite{R07, S12}, customer service \parencite{R08}, and clinical decision-making \parencite{R09}, highlighting the need for context-aware frameworks. Finally, there is a recognized need to refine or develop new theoretical models that capture the non-linear and dynamic nature of GenAI adoption, moving beyond traditional technology acceptance paradigms \parencite{S13}.

\paragraph{Formation of Hybrid Human-AI Teams}
An emerging and widely discussed research avenue in the IS field concerns the formation and functioning of hybrid human-AI teams, where humans and GenAI systems collaborate in shared tasks. Several sources highlight the need to explore collaborative dynamics, including how responsibilities, decision-making authority, and handover points are distributed between human and AI agents \parencite{R05, R07, S06}. This includes investigating symbiotic relationships that augment human intelligence rather than replace it, and understanding how to design collaboration frameworks that prevent over-reliance on AI while leveraging its strengths \parencite{R07, S12}. The organizational implications of such hybrid teams are also significant, as GenAI adoption is expected to reshape business processes, IT capabilities, and workforce structures \parencite{R03, R10}. Moreover, researchers are encouraged to examine how cultural, regulatory, and sector-specific contexts influence the integration of AI into teams \parencite{R03}. Finally, the evolving composition of teams—including the competencies required to work effectively with AI agents—presents a rich area for inquiry, calling for new models of team design and skill development \parencite{S06}.

\paragraph{Automation of Routine Tasks}
A key area of interest in GenAI research within IS is the automation of routine and repetitive tasks, which has implications across organizational and managerial levels. Several sources emphasize the potential of GenAI to support or fully automate tasks in domains such as  education, enterprise management, and service operations, often through the use of engineered prompts and workflows tailored to specific functions \parencite{R02, R04, R08}. This shift invites deeper investigation into how automation can augment human capabilities, redefine job roles, or even create new forms of work, rather than merely replacing existing ones \parencite{R05}. Strategic concerns such as risk containment, accountability, and transparency in automated processes are also critical, especially in customer-facing and regulated environments \parencite{R05, R08}. Furthermore, researchers are called to explore the implications of automation, including the extent of job replacement \parencite{R10}. Finally, the relevance of automation extends to broader contexts such as smart cities and sustainable infrastructure, where GenAI can play a role in optimizing design, construction, and operations \parencite{S15}.

\paragraph{Organizational Transformations to Incorporate GenAI}
The integration of GenAI into organizational contexts is expected to drive significant transformations in business processes, structures, and strategic models. Multiple sources emphasize the role of GenAI in revealing opportunities for process innovation and supporting process (re-)\hspace{0pt}design initiatives, particularly through automation and augmentation of decision-making and resource management \parencite{R04, R07, R08}. Researchers can also explore how GenAI can facilitate digital transformation across industries, including shifts from low- to high-value services and the emergence of new business models \parencite{R07, R08}. Moreover, GenAI’s impact on organizational capabilities, including IT infrastructure, human resource management, and knowledge systems, calls for a rethinking of traditional organizational boundaries and roles \parencite{R07, R10}. Scholars are also urged to adopt a sociotechnical perspective, examining GenAI not merely as a technical tool but exploring its boundary conditions that define its presence and impact within the context of an organization \parencite{R10}.

\paragraph{Benefits Provided by GenAI to Organizations}
Future research in IS should examine GenAI’s impact across diverse sectors such as medicine, education, tourism and e-commerce \parencite{R07, S01, S14, S17} assessing its potential benefits, e.g. the ability to personalize and augment services \parencite{S17}, as well as enhanced productivity, creativity, and service quality \parencite{R02}. Importantly, these benefits must be evaluated alongside ethical considerations and societal implications, ensuring responsible deployment that aligns with values such as social justice and sustainability \parencite{R02, S17}.

\paragraph{Risks Introduced by GenAI to Organizations}
Understanding the risks and challenges associated with GenAI adoption is essential for responsible organizational integration. Key concerns include the strategic containment of automation risks, particularly in service industries where GenAI may disrupt established marketing and operational practices \parencite{R05, R08}. In high-stakes domains like healthcare, researchers are urged to address issues of accuracy, guideline compliance, and implementation challenges, ensuring safe and effective use of AI-driven tools \parencite{R09, S03}. These risks highlight the need for robust governance frameworks and continuous evaluation of GenAI’s organizational impact.

\subsubsection{Societal Perspective [F3]}

\paragraph{Societal Impact of GenAI}
The societal impact of GenAI is a multifaceted area of inquiry in IS, with researchers increasingly called to examine its implications for equity, labor markets, global development, and ethical governance. A recurring theme across sources is the need to understand how GenAI may displace or transform jobs and crowdsourced initiatives, and what welfare consequences this may entail \parencite{R04, R05, R10}. These changes raise additional concerns about "societal stress" caused by job replacement \parencite{R10}.
Another major research direction involves the global dimension of GenAI’s impact, including its role in widening or bridging the digital divide between countries at different stages of technological development, and its influence on resource allocation across the Global North–South divide \parencite{R07}. The integration of GenAI into global IT management and outsourcing structures further highlights the importance of addressing regulatory diversity, cultural sensitivity, and language-specific capabilities to ensure inclusive and deployment \parencite{R03}. Ethical tensions also emerge in scenarios where bias mitigation efforts may reduce profitability or create competitive disadvantages, prompting calls for frameworks that balance efficiency, fairness, and social responsibility \parencite{R01, S18}.
In sectors such as healthcare and education, GenAI’s societal role is particularly pronounced. Researchers emphasize the need for empirical, interdisciplinary, and user-centered studies to validate its pact on public health systems and education systems \parencite{S01, S11, S17}. These studies should account for cultural, linguistic, and socio-economic diversity, ensuring that GenAI technologies are designed and evaluated in ways that reflect real-world complexity. Finally, scholars are encouraged to adopt a sociotechnical systems perspective, recognizing GenAI as a deeply embedded actor within societal ecosystems, whose boundaries, interactions, and ethical implications must be critically examined \parencite{R09, R10}.

\paragraph{Application to Specific Domains and Business Sectors}
Overall, future IS research should focus on developing tailored GenAI solutions that are not only technically sound but also ethically responsible and contextually relevant. A recurring issue related to GenAI contextual adaptation and domain-specific performance is the need to fine-tune generative models to meet the unique requirements of sectors such as healthcare, finance, education, marketing, e-commerce, tourism, and entertainment \parencite{R01, R02, R04, R07, R08}. Researchers are encouraged to explore how GenAI can support enterprise enterprise management, as well as process design and re-design, contributing to operational efficiency and strategic innovation \parencite{R04}. In high-stakes domains, such as healthcare and finance, the development of ethical guidelines and regulatory frameworks is essential to balance the competing demands \parencite{R01, R09}.
Healthcare, in particular, emerges as a focal point for future research, with calls for interdisciplinary collaboration, clinical validation, and standardized evaluation tools to assess GenAI’s utility in diagnostics, documentation, and patient communication \parencite{S01, S02, S03, S05, S11, S16}. In education, GenAI’s potential to enhance learning experiences and develop student skills invites further investigation into pedagogical models and responsible use practices \parencite{R08, S14}. Additionally, it is worth to examine how GenAI can transform public services, including government-led digital initiatives, while minimizing misuse in sensitive domains such as legal and healthcare services \parencite{R08}.

\paragraph{Regulations and Legal Issues Related to GenAI}
Legal and regulatory issues surrounding GenAI are becoming increasingly central to IS research, as organizations and governments grapple with the challenges of ethical deployment, data governance, and intellectual property protection. A key direction involves the development of ethical guidelines and legal frameworks tailored to sectors such as healthcare and finance, where the tension between rapid deployment, accuracy, and inclusivity creates unique regulatory demands \parencite{R01}.
Another prominent area of inquiry concerns the governance of bias and fairness, including how AI laws and data filtering techniques can be designed to reduce algorithmic discrimination while maintaining transparency and accountability \parencite{R07, S18}. The rise of GenAI also raises complex questions about intellectual property rights, prompting calls for new metrics and legal interpretations that reflect the generative nature of these technologies \parencite{R07, R08}. In the context of service industries and marketing, governments and organizations must address emerging privacy, security, and user data protection challenges, while establishing clear standards for responsible use \parencite{R08, S11}.
The applications of GenAI in areas such as healthcare and legal cases further underscore the need for comprehensive regulatory oversight, which requires building interdisciplinary collaborations among technologists, professionals, ethicists, and policymakers to ensure that GenAI tools are aligned with ethical principles and legal requirements \parencite{R10, S11, S16}.

\paragraph{Risks Introduced by GenAI to Societies}
GenAI introduces a range of societal risks that warrant close attention from IS researchers. Key concerns include its potential to disrupt industries such as tourism, e-commerce, and healthcare, and to alter human interactions and decision-making in ways that may affect creativity, productivity, and social justice \parencite{R02, R07, R08}. In sensitive domains like medicine and education, scholars emphasize the need for rigorous evaluation and standardized tools to assess GenAI’s safety and effectiveness \parencite{S01, S14}. Broader issues such as privacy, security \parencite{R10}, and the impact of automation on urban environments \parencite{S15} also require investigation to ensure GenAI supports sustainable and equitable societal development.

\paragraph{New Use Cases for GenAI Applications}
Exploring new use cases for GenAI applications is a dynamic and promising research direction in IS. Scholars are increasingly interested in how GenAI can enhance existing research methods or even propose novel ones, potentially transforming the way knowledge is produced and disseminated \parencite{R02, R06}. In design science, GenAI is seen as a tool to foster creativity in the development of new IT artifacts \parencite{R04}. This opens the door to new genres of academic publication and alternative theorizing processes, which may challenge traditional norms and encourage methodological diversity \parencite{R06}.
Beyond academia, GenAI is being applied to address global grand challenges, such as environmental protection and the Sustainable Development Goals, by expanding modes of explicit knowledge production and improving efficiency in problem-solving \parencite{R07}. In management and cybersecurity research, generative algorithms are used to simulate cyber-attacks, generate marketing content, and explore social media dynamics, demonstrating their versatility in both analytical and creative tasks \parencite{S04}. In education, GenAI supports digital pedagogical innovations, blending AI-driven assistance with traditional teaching to create modern, adaptive learning environments \parencite{S12}. In healthcare, emerging use cases include remote patient monitoring and predictive analytics, which promise to enhance care delivery and operational efficiency \parencite{S16}. All such new and innovative applications require further research to enable both efficient and safe use of GenAI.

\paragraph{Effect on Job Market}
GenAI is expected to significantly reshape the job market, particularly by automating routine tasks and transforming roles in information-intensive sectors \parencite{R04, R07}. While some jobs may be replaced, new roles will likely emerge that focus on collaborating with AI systems and leveraging human-specific skills \parencite{R07}. In consumer-facing services, GenAI may alter employee responsibilities and require reskilling to adapt to changing workplace dynamics \parencite{R08}. As the impact of GenAI on the job market can result in tensions and resistance in societies, this topic is of particular interest to researchers.

\subsubsection{Ethical Perspective [F4]}

\paragraph{Ethical Use of GenAI}
Future research on the ethical use of GenAI should aim to balance innovation with societal responsibility. Achieving this will require interdisciplinary research focusing on both systemic frameworks and user-centered practices. Key priorities include developing comprehensive ethical and regulatory guidelines to address privacy, data security, and fairness, alongside investigating user training to mitigate misuse, especially by non-technical audiences \parencite{R01,S06,S11,S12,R07}. Researchers should also empirically evaluate social, psychological, and economic impacts while exploring complementary ethical approaches, taking into account diverse stakeholder perspectives \parencite{R02,R09,S09}. The emergence of explainable GenAI (GenXAI) highlights a need for transparent, contextualized explanations and effective bias mitigation in high-stakes domains \parencite{S18}. Across sectors—from IT outsourcing to education—future work must ensure ethical integration of GenAI technologies, promoting social justice and operational accountability as commercial adoption accelerates \parencite{R03,R10,R08}.

\paragraph{GenAI Bias Mitigation}
Future research on mitigating bias in GenAI, especially large language models (LLMs), calls for multidisciplinary approaches focusing on dynamic, scalable, and context-aware methods. These should exploit various approaches, from real-time user feedback systems to resource-efficient longitudinal studies for bias detection and mitigation \parencite{R01}. 
Understanding bias requires addressing how foundational model biases propagate into downstream applications where speed and scalability often challenge fairness objectives \parencite{R01,R09}.
Dealing with the problem of bias also requires prior research on bias sources and types \parencite{R01,R02,R07,R08}. Such research could be supported with explainable AI \parencite{S18}. Understanding bias across cultural and national contexts requires balancing localized adaptations with global standards \parencite{R01}.
Ethical governance is critical for effective bias mitigation, requiring multi-stakeholder frameworks emphasizing transparency, accountability, explainability, human oversight, privacy, and the right to contest AI outcomes, all embedded in informed regulatory standards and laws tailored for GenAI \parencite{R07, S11, S17}.


\paragraph{Transparency and Explainability}
Personalization of explanations based on user expertise and context is a critical direction for enhancing accessibility and ensuring that diverse audiences could meaningfully interpret model outputs \parencite{S18, S09, R08}. 
Design science perspectives can support the development of innovative techniques that strengthen user understanding, adoption, and trust in these systems \parencite{R02, S17}. In order to improve interpretability while maintaining performance, explainable AI methods should be considered \parencite{R01, R02}. Apart from technical solutions, developing comprehensive ethical frameworks and guidelines addressing transparency should be considered \parencite{S12}.
The explainability is highlighted as a crucial factor for GenAI's healthcare applications, where future research studies should examine how explainability affects clinician reliance on automated diagnostics \parencite{R09, S01, S07, S16}. 


\paragraph{Awareness of User’s Specifics}
Research on user-specific awareness in GenAI should explore applications where GenAIs interact with diverse users in various contexts, such as employee recruitment, credit scoring, customer service, sentiment analysis, and content recommendation \parencite{R01}. The lack of such awareness can affect fairness of such GenAI solutions.
Integrating generative models within outsourcing and localization strategies demands sensitivity to local regulations, cultures, and language capabilities \parencite{R03}. Localization extends to the adaptation of GenAI tools for new languages and new communication contexts, putting stress on appropriate translation and cultural relevance \parencite{S11}.
This is particularly relevant in healthcare applications that require careful evaluation across linguistic, cultural, and socio-economic landscapes, with attention to environmental sustainability \parencite{S17}. The design of GenAI must consider whether models are built for static or dynamic contexts, further complicating the awareness issue and inducing likely fairness and transparency challenges \parencite{R07}. Finally, future work should develop explanatory frameworks attuned to user needs and ethical-social factors, as generative models mature \parencite{S18}.

\paragraph{Proper Representation and Inclusivity}
Inclusivity-oriented research is needed to mitigate data biases that result in unequal service or treatment quality and to strengthen data governance practices \parencite{R01, R09, S11}. Similarly, the perspectives of underrepresented users on AI ethics — especially regarding explainability and perceived fairness — require systematic investigation \parencite{S09}.
Scholars should examine how GenAI can enable culturally adaptive and inclusive digital localization processes that serve global audiences while preserving local authenticity \parencite{R03}. 
Future research on proper representation and inclusivity in AI faces a key challenge in developing ethical and fairness-aware frameworks for LLMs across domains such as healthcare, finance, and marketing, where rapid deployment often conflicts with equity objectives \parencite{R01}. Research should further explore how to distinguish between purposeful and unintended differentiation in such systems to prevent inequitable outcomes for marginalized groups \parencite{R01}.


\subsubsection{Engineering Perspective [F5]}

\paragraph{Definition of GenAI Metrics}
A strong message about the need for new metrics that would enable evaluation of GenAIs. Such dedicated metrics could provide a base for specialized tools to assess benefits and risks of GenAI models like ChatGPT in various fields \parencite{S14}.
Traditional measures, such as accuracy, precision, and recall are insufficient for GenAI tasks. Therefore, new metrics incorporating helpfulness, harmlessness, honesty (HHH), security, and standardized validation on independent datasets are necessary \parencite{R02, S02, S07}.
These should address the impact of GenAI on individuals, workers, and organizations, as well as the inter-organizational impacts \parencite{R10}.
Regarding individuals, measuring GenAI's influence on cognitive aspects such as questioning, rigor, and clarity is underexplored and warrants specific evaluative frameworks \parencite{R06}.
Accuracy assessment, especially concerning AI hallucinations in generative models, remains a challenge and calls for dedicated reliability metrics \parencite{R08}.
New metrics are also needed to measure fairness and diversity in LLM-driven systems \parencite{R01} and explainable AI could help in their application \parencite{S18}.
Also intellectual property rights protection demands comprehensive metrics to safeguard legal and ethical standards in GenAI applications \parencite{R07}.
Lastly, refinement of theoretical frameworks on technology adoption is needed to better capture the nuances of GenAI integration and acceptance in diverse environments \parencite{S13}. 

\paragraph{Empirical Evaluation of GenAI in the Field}
The empirical evaluation of GenAI systems calls for developing resource-efficient longitudinal designs that could monitor how GenAI tools mitigate bias over time while accounting for expertise gaps and computational constraints \parencite{R01}. 
It should definitely include user-centered aspects of GenAI systems, particularly the ethical and privacy implications of conversational models. Future studies should move beyond simulated testing to investigate real-world user interactions, collecting authentic behavioral data while safeguarding user privacy and ensuring ecological validity \parencite{S08}. 
There is also a recognized need to empirically assess whether GenAI-based tools genuinely enhance academic rigor, questioning, and clarity in scholarly inquiry \parencite{R06}. 
In the healthcare domain, clinical trials, observational studies, and cross-institutional collaborations could help determine the real clinical utility of language models \parencite{S01}. Achieving this goal requires engaging healthcare professionals directly involved in documentation and decision-making processes, as well as expanding the scope of research to encompass diverse medical contexts \parencite{S05}. A coherent approach to validating AI systems in health communication must also consider technical accuracy, patient satisfaction, and public health outcomes through real-world trials \parencite{S11}. 
GenAI-driven personalized healthcare services should be studied across cultural and socio-economic boundaries, evaluating their sustainability and impact on equitable care delivery \parencite{S17}. 

\paragraph{Design Principles for GenAI Solutions}
Research on design principles for GenAI should first address the question of effective design principles that guide GenAI development holistically, integrating human-centered and technical perspectives \parencite{R04}.  
An important research direction explores how to design and test GenAI-based systems with varying levels of automation to optimize human benefit considering individual diversity \parencite{R05}. This personalization-oriented approach complements efforts focused on defining the most effective design processes for creating GenAI that operates collaboratively with humans, balancing autonomy and human oversight \parencite{R07}.  
Another theme is the simplification of model architectures, improving data quality, and implementing standardized validation procedures in a strive to ensure that GenAI systems are reliable, transparent, and suitable for critical settings such as healthcare \parencite{S07}. 
Moreover, privacy-oriented design emerges as a necessary complement to usability and transparency research. Addressing gaps in privacy-aware design demands systematic exploration of how user-privacy principles can be embedded throughout the development lifecycle, from early prototyping to deployment. This includes not only identifying potential risks but also developing practical techniques, frameworks, and design guidelines for privacy-sensitive interfaces like chatbots \parencite{S08}.  

\paragraph{GenAI Model Training}
Future research on GenAI model training converges around several key domains: model adaptation, bias mitigation, domain specialization, and ethical application.  
One major direction involves developing strategies for effective model fine-tuning to balance accuracy and generalization, without overfitting or sacrificing scalability \parencite{R01}. 
Closely related to these concerns is the effort to fine-tune or adapt models for specific domains and contexts \parencite{R02, R04}. Research should examine whether GenAI is best suited for static or dynamic environments, and how models can evolve continuously without structural conflicts during extended use \parencite{R07}. 
The integration of GenAI in outsourcing sector necessitates to train the models considering global regulatory variations, cultural differences, and language specific capabilities while addressing efficiency and innovation \parencite{R03}.  
The challenge of preventing bias in algorithms and data used to train models persists as a critical focus area \parencite{R01,R08}. 
Healthcare is an exemplary target area for which domain-specific improvements in model training should be considered. These include enhancing model accuracy, reliability, and robustness for safe clinical deployment \parencite{S16}, guidance by up-to-date medical standards \parencite{S03}, standardized validation and diverse training datasets representing multiple disease categories for supporting reproducibility and generalization of research \parencite{S07}, as well as advanced multimodal learning integrating textual and visual medical data to improve diagnostic understanding \parencite{S11}. As an alternative (or a support) to domain-specialized models, Retrieval‑Augmented Generation could be considered \parencite{S02}.  
As regards ethical application, guidelines for the responsible integration of LLMs should be developed for high-stakes sectors such as healthcare and finance, emphasizing the balance between rapid deployment, accuracy, and inclusivity with regulatory considerations \parencite{R01}.


\paragraph{GenAI System Development and Maintenance}
A promising future research direction is the exploration of GenAI's role in design science research to enhance creativity in developing new IT artifacts which could lead to advancing the theoretical and practical foundations of IT development \parencite{R04}.
A related focus is the development and testing of partially automated tools aimed at maximizing human benefit. Specifically, advancing the technical integration of AI within robotic process automation (RPA) tools could lead to more sophisticated, adaptive automation solutions in organizational contexts \parencite{R05}.
In terms of system integration, future work should target addressing interoperability, usability, and compliance challenges. For instance, in the case of existing healthcare information systems, including electronic health records (EHRs), it would enable GenAI to augment clinical workflows and decision-making \parencite{S11}.


\subsubsection{Quality Requirements Perspective [F6]}

\paragraph{Data Privacy Protection}
Data privacy protection should be considered of primary importance in future research on GenAI due to the risks and harms, especially the ones arising from real-time interactions with chatbots. Current studies often rely on simulated data, leaving a gap for research based on actual user data in real-world settings \parencite{S08}. Future work should focus on developing user-privacy-centric chatbot designs, incorporating robust privacy safeguards and intellectual property protections \parencite{S08,R08}.
Personalization and user interaction improvements, particularly in health communication, should ensure that privacy and consent safeguards are in place to deem the related GenAI solutions trustworthy. There is a need to study GenAI telehealth applications across diverse cultural and socio-economic contexts to evaluate informed clinical decision-making and sustainability \parencite{S11,S17}.
The broader societal impact includes transformative changes in workforce and operational processes across sectors. This calls for renewed research into end-user computing and examination of societal issues like privacy and security in the pervasive use of GenAI \parencite{R03,R10}.
The growing sophistication of GenAI-powered privacy-related cyber threats requires IT outsourcing firms leveraging GenAI to develop advanced real-time detection, response, and mitigation systems to protect sensitive data and maintain trust with clients \parencite{R03}.
There is a growing demand for comprehensive guidelines addressing data privacy protection that would define accountability structures involving all stakeholders to ensure safe and transparent deployment \parencite{S11,S12,S17}.
This could be undertaken as a part of a wider effort on data governance concerning AI, especially in healthcare, which remains an understudied area in IS research \parencite{R09}.


\paragraph{Security and Protection}
Security challenges are amplified by GenAI's capabilities, leading to sophisticated cyber threats. The development of advanced AI-driven real-time detection, response, and mitigation methods is critical, alongside protocols against attacks such as prompt injections and data poisoning. Human oversight and clear accountability frameworks are essential to secure GenAI systems effectively \parencite{R03,S17}. 
This context necessitates a deeper understanding of broader societal impacts, especially concerning the security of business applications \parencite{R10}.
Furthermore, strengthening defenses against GenAI-fueled security threats is of paramount importance. As adversarial actors exploit generative tools to scale cyber attacks, new AI-driven mechanisms for real-time threat detection, adaptive response, and mitigation must be developed \parencite{R03}. This effort should align with a sociotechnical assurance framework emphasizing human oversight, resilience against data poisoning and prompt injection, and clear accountability chains \parencite{S17}.

A closely related yet distict research direction should concern combating the misuse of GenAI. 
Governments must proactively prepare for the potential misuse of GenAI by establishing comprehensive regulations and guidelines that address its deployment in sensitive services such as healthcare and legal sectors. This involves creating frameworks to minimize generative tools' misuse, ensuring robust oversight and accountability. Additionally, preventative measures should be implemented to guard against malicious uses of AI systems, fostering a responsible environment that prioritizes ethical considerations and protects individuals and institutions from harm \parencite{R08}.
Security also gives reasons to design and enforce comprehensive frameworks to guide the responsible development and deployment of GenAI systems. These frameworks should address data protection and transparency while ensuring accountability among developers and institutional actors \parencite{S11, S12}. Future studies should clarify how such ethical principles can be operationalized in high-stakes fields like healthcare and academic publishing, ensuring that GenAI tools uphold the expected integrity and respect data sovereignty. The proposal of a modern Turing test to detect AI-generated research submissions further highlights the urgency of maintaining integrity in scholarly communication \parencite{R06}.
As GenAI is increasingly used in decision systems, including the area of security, investigators should also strive to address biases in GenAI-powered fraud detection systems which can result in higher false positive rates for transactions from certain demographic groups, leading to discriminatory practices \parencite{R01}.


\paragraph{Demonstrating Trustworthiness}
Transparency is considered fundamental for demonstrating GenAI's trustworthiness, especially in healthcare where the need for interpretable models that clearly communicate decision-making processes is crucial to foster clinician acceptance \parencite{R02,S01,S16}. 
Explainability mechanisms are equally necessary for building confidence among users, especially healthcare professionals, particularly as design science researchers develop new methods to enhance these capabilities and study their effects on adoption \parencite{R02,R04}.
Reliability and accuracy improvements cannot be ignored either. Studies should evaluate GenAI's performance across broader domains, e.g. medical specialties which require complex diagnostic and treatment decisions where accuracy is extremely important \parencite{S02,S16}. This could be supported with simplifying model architectures, implementing standardized validation procedures, and addressing data quality challenges to ensure models meet clinical standards \parencite{S05,S07}.
Mitigating hallucinations and inaccuracies in LLM outputs demands urgent attention as these issues can undermine trust \parencite{R02, S18}. Researchers should strive to develop techniques beyond Chain-of-Thought reasoning to mitigate inaccuracies and improve verifiability \parencite{S18, R08}.

Human-AI interaction research should investigate user attitudes toward GenAI, examining why people appreciate or avoid these tools and exploring trust-related factors \parencite{R02,R09}. This includes studying impacts on business workers and general users, requiring renewed focus on end-user computing topics in the GenAI era \parencite{R10}. Design approaches that foster trust through improved system reliability and transparency will be essential \parencite{R04}. Future work should also address bias issues affecting GenAI integration e.g. in clinical workflows \parencite{R09}.
Understanding the factors that influence user trust in LLM outputs—despite potential hallucinations and training data inaccuracies—is essential \parencite{R02}.


\paragraph{Efficiency and Scalability}
Future research on GenAI efficiency and scalability encompasses several interconnected dimensions. Among them, operational integration emerges as a critical theme, focusing on how GenAI can transform digital services through enhanced efficiency and operational capabilities \parencite{R08}. This includes investigating government adoption strategies for building scalable digital services accessible to broader populations \parencite{R08}. Parallel to service delivery, the localization domain demonstrates significant potential, where GenAI integration enhances efficiency, scalability, and cultural customization, enabling hyper-localized content creation \parencite{R03}.
Quality assurance and reliability represent another essential research direction, particularly in high-stakes applications. Prioritizing model accuracy, reliability, and robustness ensures safe and effective clinical applications \parencite{S16}. This intersects well with bias mitigation challenges, prompting researchers to design dynamic methodologies that integrate real-time user feedback for continuous detection and mitigation of LLM bias while maintaining scalability and operational efficiency \parencite{R01}.
Beyond individual organizational contexts, inter-organizational dynamics call for systematic investigation. As GenAI permeates organizations, interaction patterns between and among entities will evolve, necessitating research into efficient cooperation across organizational boundaries \parencite{R10}.

\paragraph{Accountability and Contestability}
A valid research challenge lies in balancing competing priorities: high-speed decision-making versus ethical accountability, and hyper-personalization versus inclusivity \parencite{R01}. In this context, operationalizing bias mitigation is important yet requires transparent, explainable interfaces that preserve human agency while implementing robust safeguards \parencite{S17}.
Research should clarify the obligations of GenAI providers to determine which GenAI applications should face restrictions or prohibition \parencite{R07}. This involves establishing clear accountability frameworks that delineate responsibilities among developers, providers, and regulators \parencite{S17}. Another identified research direction is addressing the challenge of responsibly automating organizations' business processes while ensuring transparency and accountability in GenAI deployment \parencite{R08}.
Future work should also advance reliability through information resilience protocols, including security measures against prompt injection and data poisoning \parencite{S17}. Essential components include human-in-the-loop oversight and contestability mechanisms that enable stakeholders to challenge AI-driven outcomes \parencite{S17}.

